% Related work. Only works listed in refs/citation_ledger.csv are cited.

\section{Related Work}
\label{sec:related}

\textbf{WSVAD and CLIP-based detectors.} Most WSVAD methods treat a video as a bag of snippets and train a scorer with a multiple-instance objective over the highest-scoring snippets~\cite{sultani2018real}. Later work refines the snippet representation through feature-magnitude learning~\cite{tian2021rtfm}, memory units with uncertainty regulation~\cite{zhou2023urdmu}, and similarity and distance graphs for temporal context~\cite{wu2020xdviolence}. With CLIP~\cite{radford2021clip} features, CLIP-TSA adds temporal self-attention~\cite{joo2023cliptsa}, and VadCLIP~\cite{wu2024vadclip} adds a semantic output that matches snippets against learnable class prompts. Follow-up work enriches the prompts with normality guidance~\cite{yang2024tpwng} or spatio-temporal cues~\cite{wu2024stprompt}. These methods all model time after the fixed-length resampling of snippet features. We change what reaches that step and which parameters a semantic objective may update.

\textbf{Temporal modelling and lightweight adaptation.} Explicit temporal differences are an established cue in action recognition, where TDN computes frame differences inside a trainable backbone~\cite{wang2021tdn} and TSM shifts channels across time~\cite{lin2019tsm}. In WSVAD the backbone is frozen, and the only temporal operation before learning is aggregation; our refinement targets that step. Parameter-efficient methods insert small modules or low-rank updates into a frozen model~\cite{houlsby2019adapter,hu2021lora}, and CLIP-Adapter~\cite{gao2024clipadapter} applies the idea to CLIP. Instead of adding a module, our second design restricts an extra objective to an existing adapter selected by its position in the prediction pathways.
